% !TEX root = main.tex
\documentclass{scrartcl}
\usepackage{fontspec} % Enable full UTF-8 support
\usepackage{indentfirst}

\KOMAoptions{
	% DIV=20,
	paper=a4,
	fontsize=11pt,
	%TODO: rimettere a true
	twoside=false,
	%decide se lasciare spazio in vertiale o in orizzontale nei paragrafi, in questo modo indenta la prima riga
	parskip=false,
	bibliography=totoc,
	listof=totoc,
	listof=entryprefix,
	titlepage=false,
	captions=tableabove,
	headings=standardclasses,
	% draft=true
}
\addtokomafont{title}{\normalfont\rmfamily}
\addtokomafont{subject}{\normalfont\rmfamily}

\PassOptionsToPackage{hyphens}{url}
\usepackage[a-3b,pdf17]{pdfx}
\hypersetup{hidelinks}

\begin{document}

%% TITOLO
\title{\LARGE\textbf{\uppercase{TODO: titolo della tesi}}} % TODO
\author{
	\begin{tabular}{@{}c@{}}
		TODO: Nome Cognome (matricola: TODO) % TODO
	\end{tabular}
	\\\\\vspace{1em}
	\begin{tabular}{@{}l@{}c@{}r@{}}
		Relatore             & \hspace{10em} & Correlatore           \\
		Prof.~TODO: Relatore & \hspace{10em} & Dr.~TODO: Correlatore \\ % TODO
	\end{tabular}
}
\date{}
\maketitle

%%%%CONTENUTI%%%%
% Struttura suggerita per il riassunto (~1 pagina A4):
% 1. Contesto e motivazione (1-2 paragrafi)
% 2. Sfide e gaps nella letteratura (1 paragrafo)
% 3. Lavoro svolto — articolato con \begin{enumerate}\item...\end{enumerate}
% 4. Risultati ottenuti (1 paragrafo)
% 5. Sviluppi futuri (facoltativo, 2-3 righe)

% TODO: inserire contenuto


\end{document}